% concatenated from neurips_2026_arxiv.tex
\documentclass{article}

% if you need to pass options to natbib, use, e.g.:
\PassOptionsToPackage{numbers, compress}{natbib}
% before loading neurips_2026

% The authors should use one of these tracks.
% Before accepting by the NeurIPS conference, select one of the options below.
% 0. "default" for submission
% \usepackage{neurips_2026}
% the "default" option is equal to the "main" option, which is used for the Main Track with double-blind reviewing.
% 1. "main" option is used for the Main Track
%  \usepackage[main]{neurips_2026}
% 2. "position" option is used for the Position Paper Track
%  \usepackage[position]{neurips_2026}
% 3. "eandd" option is used for the Evaluations & Datasets Track
 % \usepackage[eandd]{neurips_2026}
% 4. "creativeai" option is used for the Creative AI Track
%  \usepackage[creativeai]{neurips_2026}
% 5. "sglblindworkshop" option is used for the Workshop with single-blind reviewing
 % \usepackage[sglblindworkshop]{neurips_2026}
% 6. "dblblindworkshop" option is used for the Workshop with double-blind reviewing
%  \usepackage[dblblindworkshop]{neurips_2026}

% After being accepted, the authors should add "final" behind the track to compile a camera-ready version.
% 1. Main Track
% \usepackage[main, final]{neurips_2026}
% 2. Position Paper Track
%  \usepackage[position, final]{neurips_2026}
% 3. Evaluations & Datasets Track
 % \usepackage[eandd, final]{neurips_2026}
% 4. Creative AI Track
%  \usepackage[creativeai, final]{neurips_2026}
% 5. Workshop with single-blind reviewing
%  \usepackage[sglblindworkshop, final]{neurips_2026}
% 6. Workshop with double-blind reviewing
%  \usepackage[dblblindworkshop, final]{neurips_2026}
% Note. For the workshop paper template, both \title{} and \workshoptitle{} are required, with the former indicating the paper title shown in the title and the latter indicating the workshop title displayed in the footnote.
% For workshops (5., 6.), the authors should add the name of the workshop, "\workshoptitle" command is used to set the workshop title.
% \workshoptitle{WORKSHOP TITLE}

% "preprint" option is used for arXiv or other preprint submissions
\usepackage[preprint]{neurips_2026}

% to avoid loading the natbib package, add option nonatbib:
%    \usepackage[nonatbib]{neurips_2026}

\usepackage[utf8]{inputenc} % allow utf-8 input
\usepackage[T1]{fontenc}    % use 8-bit T1 fonts
\usepackage{hyperref}       % hyperlinks
\usepackage{url}            % simple URL typesetting
\usepackage{booktabs}       % professional-quality tables
\usepackage{amsfonts}       % blackboard math symbols
\usepackage{nicefrac}       % compact symbols for 1/2, etc.
\usepackage{microtype}      % microtypography
\usepackage{xcolor}         % colors


% my packages
\usepackage{amsmath,amssymb,bm,mathtools}
\usepackage{amsthm}
\usepackage{multirow}
\usepackage{graphicx}

\graphicspath{{fig/}}
\providecommand{\E}{\mathbb{E}}
\providecommand{\I}{I}
\newcommand{\norm}[1]{\left\lVert #1 \right\rVert_2}

\theoremstyle{plain}
\newtheorem{proposition}{Proposition}[section]
\theoremstyle{remark}
\newtheorem{remark}[proposition]{Remark}



% Note. For the workshop paper template, both \title{} and \workshoptitle{} are required, with the former indicating the paper title shown in the title and the latter indicating the workshop title displayed in the footnote. 
\title{Neural Preconditioned Born Series: A Metric-Matched Framework for Learning-based Preconditioners}


% The \author macro works with any number of authors. There are two commands
% used to separate the names and addresses of multiple authors: \And and \AND.
%
% Using \And between authors leaves it to LaTeX to determine where to break the
% lines. Using \AND forces a line break at that point. So, if LaTeX puts 3 of 4
% authors names on the first line, and the last on the second line, try using
% \AND instead of \And before the third author name.


% \author{
%   Juntao Wang \\
%   School of Mathematics, Jilin University, Changchun 130012, China \\
%   Shenzhen Loop Area Institute, Shenzhen 518038, China \\
%   \And
%   Jiwei Jia\thanks{Corresponding authors: \texttt{jiajiwei@jlu.edu.cn} (Jiwei Jia); \texttt{xinliang.liu@ouc.edu.cn} (Xinliang Liu).} \\
%   School of Mathematics, Jilin University, Changchun 130012, China \\
%   Shenzhen Loop Area Institute, Shenzhen 518038, China \\
%   \And
%   Xinliang Liu\textsuperscript{*} \\
%   Ocean University of China, Qingdao 266100, China
% }

\author{
  Juntao Wang \\
  School of Mathematics, Jilin University\\
  Changchun 130012, China \\
  Shenzhen Loop Area Institute\\ 
  Shenzhen 518038, China \\
  \And
  Jiwei Jia\thanks{Corresponding authors: \texttt{jiajiwei@jlu.edu.cn} (Jiwei Jia); \texttt{xinliang.liu@ouc.edu.cn} (Xinliang Liu).} \\
  School of Mathematics, Jilin University\\ 
  Changchun 130012, China \\
  Shenzhen Loop Area Institute\\ 
  Shenzhen 518038, China \\
  \And
  Xinliang Liu\textsuperscript{*} \\
  Ocean University of China\\ 
  Qingdao 266100, China
}

% \author{
%   Juntao Wang\textsuperscript{a,c},
%   Jiwei Jia\textsuperscript{a,c}\thanks{Corresponding authors: \texttt{jiajiwei@jlu.edu.cn} (Jiwei Jia); \texttt{xinliang.liu@ouc.edu.cn} (Xinliang Liu).},
%   Xinliang Liu\textsuperscript{b}\footnotemark[1]\\[2pt]
%   \textsuperscript{a} School of Mathematics, Jilin University, Changchun 130012, China\\
%   \textsuperscript{b} Ocean University of China, Qingdao 266100, China\\
%   \textsuperscript{c} Shenzhen Loop Area Institute, Shenzhen 518038, China\\
% }
% \author{
%   Juntao Wang\textsuperscript{$\heartsuit \spadesuit$},
%   Jiwei Jia\textsuperscript{$\heartsuit \spadesuit$}\thanks{Corresponding authors: \texttt{jiajiwei@jlu.edu.cn} (Jiwei Jia); \texttt{xinliang.liu@ouc.edu.cn} (Xinliang Liu).},
%   Xinliang Liu\textsuperscript{$\clubsuit$}\footnotemark[1]\\[2pt]
%   \textsuperscript{$\heartsuit$} School of Mathematics, Jilin University, Changchun 130012, China\\
%   \textsuperscript{$\clubsuit$} Ocean University of China, Qingdao 266100, China\\
%   \textsuperscript{$\spadesuit$} Shenzhen Loop Area Institute, Shenzhen 518038, China\\
% }

\begin{document}


\maketitle


\begin{abstract}
High-frequency Helmholtz problems in heterogeneous media remain challenging for both classical iterative methods and end-to-end neural PDE solvers.
We propose Neural Preconditioned Born Series (NPBS), a learned iterative preconditioning framework that operates in preconditioned residual coordinates induced by the Convergent Born Series (CBS).
Existing learned Born-series methods primarily use Born-style unrolling for forward wavefield prediction, while learned Helmholtz preconditioners are usually formulated in physical residual coordinates.
NPBS fills this gap by recasting Born-series iteration as shifted-Laplacian left preconditioning, and replacing the CBS preconditioner with a learned residual-to-correction map in the Born-preconditioned coordinates.
The left preconditioner further induces a residual metric, which yields a metric-matched training objective that aligns optimization with the preconditioned geometry used at inference.
On heterogeneous Helmholtz benchmarks, metric-matched NPBS reduces iteration counts by up to $1.9\times$ over direct residual learning, with gains increasing from $1.2\times$ to $1.9\times$ as the wavenumber rises.
Compared to classical CBS, learned NPBS reduces stationary iteration counts by over $20\times$; when used as a preconditioner for FGMRES, it further achieves the lowest wall-clock time among all evaluated methods.
The same metric-matched formulation also improves convergence on convection--diffusion--reaction systems and Newton linear systems for nonlinear PDEs, indicating that residual-metric matching is a general design principle for neural preconditioners.
\end{abstract}

\section{Introduction}
Efficient solution of wave equations in heterogeneous media is central to a broad range of engineering and scientific applications, from seismic inversion and medical ultrasound imaging to electromagnetic scattering and photonic design. Many of these problems reduce to the heterogeneous Helmholtz equation, whose high-wavenumber and high-contrast regimes remain challenging for both classical and learning-based solvers~\cite{gander2019class}.
The main obstacles are well known: standard discretizations suffer from pollution errors~\cite{babuska1997pollution}, multilevel methods face coarse-grid dispersion mismatch~\cite{stolk2016dispersion}, and indefiniteness can undermine relaxation and coarse correction~\cite{ernst2011difficult}. Classical approaches address these issues through shifted-Laplacian preconditioning~\cite{erlangga2006novel,erlangga2004class}, wave-ray multigrid~\cite{brandt1997wave}, Convergent Born Series (CBS)~\cite{osnabrugge2016convergent}, and domain decomposition~\cite{chen2013source,graham2020domain}. Despite these advances, robust and efficient solution of heterogeneous Helmholtz problems remains difficult.

Learning-based PDE solvers have developed rapidly in recent years. Physics-informed neural networks (PINNs)~\cite{raissi2019physics} encode governing equations and boundary conditions directly into the training objective, but often require a separate optimization for each new instance, limiting cross-instance generalization and amortized efficiency~\cite{sun2020surrogate}. Neural operators address this limitation by learning mappings between function spaces, allowing one model to solve families of PDEs parameterized by coefficients, sources, or geometries. Representative works include DeepONet~\cite{lu2021learning}, FNO~\cite{li2020fourier,li2023fourier}, UNO~\cite{rahman2022u}, and MgNO~\cite{he2024mgno}, among many others~\cite{kovachki2023neural,cao2021choose,MELCHERS2026118893}.
While neural operators greatly improve scalability, end-to-end operator learning remains difficult for highly oscillatory solutions. Because of the spectral bias of neural networks~\cite{rahaman2019spectral,liu2024mitigating}, high-frequency components are typically much harder to learn accurately than low-frequency ones. For high-frequency Helmholtz equations, this makes one-shot wavefield prediction particularly challenging~\cite{he2025self}.

A promising alternative is to embed a learned module inside an iterative solver, using the network as a residual-to-correction map rather than as a direct solution predictor~\cite{kopanivcakova2025leveraging,rudikov2024neural,zhang2024blending}. Recent learned iterative solvers have shown that this strategy can reduce iteration counts or wall-clock time relative to classical baselines~\cite{xie2025mgcfnn}, making this direction especially promising for challenging Helmholtz problems.
Our work follows this direction, but asks a more specific question: what residual coordinates and training metric should a neural preconditioner use for ill-conditioned Helmholtz problems?

We answer this question by proposing \emph{Neural Preconditioned Born Series} (NPBS), a learned iterative preconditioning framework inspired by the CBS.
Our key observation is that the Born series is algebraically equivalent to shifted-Laplacian left preconditioning; in particular, the Born residual coincides with the residual of the left-preconditioned system.
This equivalence identifies the Born-preconditioned residual coordinates as a natural state space for learned iterative correction.
NPBS builds on this view by replacing the CBS preconditioner with a neural residual-to-correction map acting on Born-preconditioned residuals.
The same equivalence also determines how the model should be trained.
If inference is performed in Born-preconditioned coordinates, then the training loss should measure residual error in the geometry induced by the same preconditioner, rather than in the unpreconditioned Euclidean metric.
Let $G_\eta$ denote the Green operator of the damped reference operator.
We introduce the residual metric induced by $G_\eta$:
\begin{equation}
\|\bm r\|_{R_\eta} := \|G_\eta \bm r\|_2,
\qquad
R_\eta := G_\eta^\ast G_\eta,
\end{equation}
and derive a metric-matched Born-series objective $\mathcal{L}_{\mathrm{bs}}^{R_\eta}$ that aligns training and inference within the same preconditioned geometry.
This alignment is especially important for indefinite Helmholtz problems, where different modes can be amplified highly nonuniformly near resonance~\cite{ernst2011difficult,erlangga2004class}.
Although developed first for Helmholtz equations, NPBS also extends naturally to convection--diffusion--reaction (CDR) equations and Newton linear systems for nonlinear PDEs~\cite{hao2024newton,lee2025neural}.


Our contributions are summarized as follows:
\begin{enumerate}
\item \textbf{A unified learned iterative framework in Born-preconditioned coordinates.}
We introduce NPBS, which lifts the diagonal CBS preconditioner to an architecture-agnostic learned residual-to-correction map operating on Born-preconditioned residuals.

\item \textbf{A metric-matched training principle induced by the reference Green operator.}
By identifying Born residuals with shifted-Laplacian left-preconditioned residuals, we derive the induced metric $R_\eta = G_\eta^\ast G_\eta$ and the objective $\mathcal{L}_{\mathrm{bs}}^{R_\eta}$, making residual-metric matching the principle that aligns training with preconditioned inference.

\item \textbf{Robustness to ill-conditioning and algorithmic generalization.}
On heterogeneous Helmholtz benchmarks, metric-matched NPBS reduces iteration counts by up to $1.9\times$ over direct residual learning, with gains growing from $1.2\times$ to $1.9\times$ as the wavenumber rises, at only $0.7\%$ additional training cost. The advantage persists across four neural backbones; when used as a preconditioner inside FGMRES, NPBS achieves the lowest average wall-clock time among all evaluated methods. The same design principle also improves convergence on CDR equations and Newton linear systems for nonlinear PDEs.
\end{enumerate}

The paper is organized as follows. Section~\ref{sec:related_work} positions NPBS relative to neural preconditioners and Born-inspired neural solvers. Section~\ref{sec:born_identity} establishes the algebraic equivalence between Born-series residuals and shifted-Laplacian left preconditioning. Section~\ref{sec:npbs} introduces the NPBS iteration and derives the metric-matched training objective. Section~\ref{sec:numerical_results} reports experiments on heterogeneous Helmholtz benchmarks across difficulty levels and backbone architectures, with extensions to CDR and nonlinear PDEs.

\section{Related Work}
\label{sec:related_work}

\paragraph{Neural preconditioners and learned iterative solvers.}
Recent learning-based solvers increasingly use neural networks as components of an iterative numerical method rather than as standalone solution predictors. Examples include learned multigrid and meta-multigrid schemes~\cite{greenfeld2019learning,chen2022meta}, neural-operator-accelerated Krylov methods~\cite{rudikov2024neural,kopanivcakova2025leveraging}, hybrids of neural operators and relaxation methods~\cite{zhang2024blending,lee2025hybrid}, and neural preconditioned Newton-type solvers~\cite{lee2025neural}. For Helmholtz equations, learned preconditioning has been studied through multigrid-augmented deep preconditioners~\cite{azulay2022multigrid,lerer2024multigrid}, Wave-ADR-NS~\cite{cui2025neural}, MGCFNN~\cite{xie2025mgcfnn}, and many others~\cite{stanziola2021helmholtz,lee2025fast}. These methods demonstrate that learned solver components can reduce iteration counts or improve amortized solve time, but they typically formulate the learned correction in the physical residual coordinates of $A\bm u=\bm f$ or focus on architecture design for approximating the inverse/preconditioner. 
In contrast, NPBS aligns its training metric with the preconditioned coordinates used at inference, rather than designing a new architecture or optimizing in unpreconditioned coordinates. This makes the formulation backbone-agnostic: the same metric-matched objective (Section~\ref{sec:npbs}) applies regardless of the correction network.

\paragraph{Born-inspired neural solvers.}
The classical Born series/Lippmann--Schwinger formulations express wave scattering through a reference Green operator and a medium perturbation; CBS~\cite{osnabrugge2016convergent} modifies this construction to obtain a convergent fixed-point method for heterogeneous Helmholtz problems. Recent learned Born-series methods, such as Neural Born Series Operator (NBSO)~\cite{zeng2023neural} and Learned Born Series (LBS)~\cite{stanziola2024learned}, use Born-style unrolling to build forward wavefield simulators for scattering media. Their central use of the Born structure is therefore as a learned finite-depth surrogate for the wavefield map. NPBS reuses the same Born-series structure not as a finite-depth forward surrogate, but as the preconditioned coordinate system for a convergent iterative solver (Section~\ref{sec:npbs}).
Consequently, NPBS can converge to the PDE solution under a verifiable residual criterion and can be embedded as a preconditioner within Krylov subspace methods (e.g.\ FGMRES), rather than serving only as a fixed-depth forward predictor.


\section{From Born Series to Preconditioned Coordinates}
\label{sec:born_identity}
%%%%%%%%%%%%%%%%%%%%%%%%%%%%%%%%%%%%%%%%%%%%%%%
\subsection{Operator Splitting and Integral Equation}
We consider the Helmholtz equation for wave scattering in an unbounded heterogeneous medium, subject to the Sommerfeld radiation condition:
\begin{equation}\label{eq:helm}
\begin{aligned}
    \mathcal{L}u(x) := -\Delta u(x) - k(x)^2 u(x) &= f(x), \quad x \in \mathbb{R}^{d},\\
    \lim_{|x| \to \infty} |x|^{\frac{d-1}{2}} \left( \frac{\partial u}{\partial |x|} - i k_0 u \right) &= 0,
\end{aligned}
\end{equation}
where $u$ is the wavefield, $f$ the source, $d$ the spatial dimension, and $k(x)=\omega/c(x)$ the spatially varying wavenumber ($\omega$: angular frequency, $c(x)$: speed of sound). The constant $k_0$ is the background wavenumber ($k(x)\to k_0$ as $|x|\to\infty$).

For computation, the unbounded domain is truncated to a bounded domain $\Omega$, and absorbing boundary layers---sponge layers~\cite{israeli1981approximation} or PML~\cite{berenger1994perfectly,stanziola2021helmholtz}---are incorporated so that the resulting bounded-domain operator (still denoted $\mathcal{L}$) suppresses spurious reflections. After discretization (finite differences or FFT-based pseudospectral methods), we obtain the linear system
\begin{equation}
A \bm u = \bm f,
\end{equation}
where $A$ is the discrete Helmholtz operator (including absorbing-layer terms), $\bm u$ the discrete unknown, and $\bm f$ the discrete source.
Introduce a damped reference operator
\begin{equation}
\mathcal{L}_\eta := -\Delta - (k_0^2 + i\eta), \qquad \eta>0,
\end{equation}
used only for preconditioning/training coordinates. The reference is chosen to be transform-diagonalizable (periodic FFT or Dirichlet/Neumann DST/DCT) for fast application; any mismatch between this reference and the physical operator, including absorbing-layer terms, is absorbed into the perturbation
\begin{equation}
\mathcal{V}_\eta := \mathcal{L}_\eta - \mathcal{L}.
\end{equation}

Let $g_\eta$ be the Green kernel for the \emph{continuous} reference operator $\mathcal{L}_\eta$ under the chosen boundary condition on $\Omega$. Define the associated continuous Green operator
\begin{equation}
(\mathcal{G}_\eta q)(x) := \int_{\Omega} g_\eta(x,y)\,q(y)\,\mathrm{d}y,
\qquad q\in L^2(\Omega),
\label{eq:G_conv}
\end{equation}
so that $\mathcal{G}_\eta = \mathcal{L}_\eta^{-1}$ on the admissible function class.

\paragraph{Discrete implementation.}
In the discrete setting, let $L_\eta$ denote the matrix obtained by discretizing $\mathcal{L}_\eta$ (e.g., via finite differences or Fourier pseudospectral methods; see Appendix~\ref{app:fft_ps} for details), and define the discrete Green operator $G_\eta := L_\eta^{-1}$. For transform-diagonalizable references (periodic FFT or Dirichlet/Neumann DST/DCT), $G_\eta$ can be applied in $\mathcal{O}(N\log N)$ operations (Appendix~\ref{app:fft_verify_all}).

\paragraph{Operator splitting and Lippmann--Schwinger form.}
With $\mathcal{V}_\eta := \mathcal{L}_\eta - \mathcal{L}$, the heterogeneous Helmholtz operator splits as
\begin{equation}
\mathcal{L} = \mathcal{L}_\eta - \mathcal{V}_\eta.
\label{eq:A_split}
\end{equation}
For the coefficient-only Helmholtz split, $\mathcal{V}_\eta$ reduces to multiplication by $k(x)^2-k_0^2-i\eta$.

In the continuous setting, $\mathcal{L}u=f$ implies
\begin{equation}
\mathcal{L}_\eta u = \mathcal{V}_\eta u + f,
\qquad
u = \mathcal{G}_\eta(\mathcal{V}_\eta u + f),
\end{equation}
equivalently,
\begin{equation}
(\mathcal{I} - \mathcal{G}_\eta \mathcal{V}_\eta)u = \mathcal{G}_\eta f.
\end{equation}
In the discrete setting, starting from $A\bm u=\bm f$ and the discrete split $A=L_\eta - V_\eta$, we obtain
\begin{equation}
(I - G_\eta V_\eta)\bm u = G_\eta \bm f.
\label{eq:LS_discrete}
\end{equation}
Equation \eqref{eq:LS_discrete} is the preconditioned operator equation used by CBS~\cite{osnabrugge2016convergent}, whose fixed-point iteration is
\begin{equation}
\bm u^{m+1}
=
\bm u^m + \gamma\!\left(G_\eta(V_\eta \bm u^m + \bm f)-\bm u^m\right),
\label{eq:CBS}
\end{equation}
where $\gamma$ is a spatially varying diagonal preconditioner derived from the perturbation $V_\eta$.

\subsection{Equivalence to Shifted-Laplacian Left Preconditioning}
We state the equivalence in discrete notation and use it to define the preconditioned learning coordinates and the induced residual metric. The accompanying spectral bounds offer quantitative intuition for why the preconditioned coordinates are favorable.

\begin{proposition}[Left-preconditioned equivalence and spectral bounds]
\label{prop:equivalence_spectral}
Let the discrete Helmholtz operator admit the split $A=L_\eta - V_\eta$, and assume $L_\eta$ is invertible so that $G_\eta := L_\eta^{-1}$ is well-defined.
Then
\begin{equation}
I - G_\eta V_\eta = G_\eta A = L_\eta^{-1}A.
\label{eq:key_identity}
\end{equation}
Hence the Lippmann--Schwinger equation \eqref{eq:LS_discrete} is algebraically identical to the preconditioned system
\begin{equation}
L_\eta^{-1}A\,\bm u = L_\eta^{-1}\bm f.
\end{equation}
Moreover, if $\|G_\eta V_\eta\|_2\le q<1$, then
\begin{equation}
\sigma(I-G_\eta V_\eta)\subset\{z\in\mathbb{C}:|z-1|\le q\},
\label{eq:disk_bound}
\end{equation}
and
\begin{equation}
\kappa_2(I-G_\eta V_\eta)\le \frac{1+q}{1-q}.
\label{eq:kappa_bound}
\end{equation}
\end{proposition}

The proof is given in Appendix~\ref{app:proof_equivalence_spectral}.

The identity \eqref{eq:key_identity} places the integral-equation formulation and the shifted-Laplacian preconditioning viewpoint in the same algebraic framework. This equivalence is exploited in two ways below: it supplies the preconditioned coordinates in which NPBS iterates (Section~\ref{sec:npbs_iter}), and it motivates the residual metric $R_\eta = G_\eta^\ast G_\eta$ used in the training loss $\mathcal{L}_{\mathrm{bs}}^{R_\eta}$ (Section~\ref{sec:bs_loss}). The spectral bounds \eqref{eq:disk_bound}--\eqref{eq:kappa_bound} quantify the benefit: whereas $A$ is indefinite, the preconditioned operator $G_\eta A$ has spectrum clustered in a disk of radius $q$ centered at $1$ and condition number at most $(1+q)/(1-q)$, reducing the complexity of the inverse map the network must learn. The formal bounds assume $q<1$; in practice, shifted-Laplacian preconditioning improves spectral concentration even beyond this regime~\cite{erlangga2006novel}, and the neural iteration exploits this structure without relying on the bound for convergence.
%%%%%%%%%%%%%%%%%%%%%%%%%%%%%%%%%%%%%%%%%%%%%%%


%%%%%%%%%%%%%%%%%%%%%%%%%%%%%%%%%%%%%%%%%%%%%
\section{Neural Preconditioned Born Series Framework}
\label{sec:npbs}

A common learned linear-solver baseline uses a neural correction map $\mathcal{M}_\theta$ driven by the \emph{unpreconditioned} residual. (Since the learned map depends on the PDE coefficients $a$, it should formally be written $\mathcal{M}_\theta(\,\cdot\,;a)$; for brevity we write $\mathcal{M}_\theta(\,\cdot\,)$ below. The specific instantiation is given in Appendix~\ref{app:architectures}.)
Given the discrete system $A\bm u=\bm f$, define the residual $\bm r^m := \bm f - A\bm u^m$ and iterate
\begin{equation}
\bm u^{m+1} = \bm u^m + \mathcal{M}_\theta(\bm r^m)
= \bm u^m + \mathcal{M}_\theta(\bm f-A\bm u^m),
\label{eq:direct_iter}
\end{equation}
which operates in Euclidean residual coordinates.

A standard training objective for $\mathcal{M}_\theta$ is to minimize the relative residual of the learned inverse action on random probes. Let $\mathcal{D}$ denote a probe distribution (e.g., $\mathcal{N}(0,I)$ or residual-replay samples). Then:
\begin{equation}
\mathcal{L}_{\mathrm{dir}}(\theta)
=
\E_{\bm r\sim\mathcal{D}}
\frac{\|A\,\mathcal{M}_\theta(\bm r)-\bm r\|_2}{\|\bm r\|_2}.
\label{eq:L_dir}
\end{equation}
This objective is simple and widely used~\cite{cui2025neural,xie2025mgcfnn}, but it optimizes the residual in \emph{unpreconditioned} Euclidean geometry; for indefinite Helmholtz operators, near-resonant modes can yield stiff gradients and slow learning.

\subsection{Neural Preconditioned Born Series (NPBS) iteration}
\label{sec:npbs_iter}

Proposition~\ref{prop:equivalence_spectral} shows that the Lippmann--Schwinger/CBS operator corresponds to shifted-Laplacian left preconditioning.
CBS uses a diagonal (pointwise) preconditioner $\gamma$ in these preconditioned coordinates; NPBS replaces $\gamma$ with a learned operator $\mathcal{M}_\theta$ acting in the \emph{same} coordinates.

\paragraph{Efficient integral-form implementation.}
\label{par:efficient_integral}
Define the Born fixed-point map $\bm u \mapsto G_\eta(V_\eta \bm u + \bm f)$.
At iterate $u^m$, NPBS updates via
\begin{equation}
\bm u^{m+1}
=
\bm u^m + \mathcal{M}_\theta\!\left(G_\eta(V_\eta \bm u^m + \bm f)-\bm u^m\right).
\label{eq:npbs_integral}
\end{equation}
Here $V_\eta \bm u^m$ is pointwise multiplication and $G_\eta$ is applied once (e.g., via FFT for transform-diagonalizable references), so each iteration requires only a single transform pair for the preconditioned residual.
If $\mathcal{M}_\theta=\gamma$, \eqref{eq:npbs_integral} reduces to CBS (\eqref{eq:CBS}).
If $\mathcal{M}_\theta = (I-G_\eta V_\eta)^{-1}$, then one NPBS step recovers the exact Helmholtz solution.

\paragraph{Equivalent left-preconditioned form.}
Using the key identity \eqref{eq:key_identity}, \eqref{eq:npbs_integral} is equivalently
\begin{equation}
\bm u^{m+1}=\bm u^m+\mathcal{M}_\theta\!\left(G_\eta(\bm f-A\bm u^m)\right),
\label{eq:npbs_iter}
\end{equation}
which makes the left-preconditioned interpretation explicit: the network input is
$\bm r_{\mathrm{bs}}^m := G_\eta(\bm f-A\bm u^m)=L_\eta^{-1}(\bm f-A\bm u^m)$.
Computing $A\bm u^m$ may require an additional application of the discrete Laplacian (e.g., an extra FFT in pseudospectral settings), hence we implement the efficient integral form \eqref{eq:npbs_integral} in practice. Proposition~\ref{prop:equivalence_spectral} implies that the preconditioned coordinates enjoy spectral clustering when $\|G_\eta V_\eta\|<1$, justifying their use for both iteration and training.

\subsection{Metric-matched training within the NPBS framework}
\label{sec:bs_loss}

We begin by analyzing the operator targets implied by the training objectives. The direct loss \eqref{eq:L_dir} trains $\mathcal{M}_\theta$ to approximate $A^{-1}$ in unpreconditioned coordinates. In contrast, the NPBS iteration \eqref{eq:npbs_integral} evaluates the network in preconditioned coordinates. Here, the effective target is $\mathcal{M}_\theta\approx (I-G_\eta V_\eta)^{-1}$ (or equivalently, $\mathcal{M}_\theta\approx (G_\eta A)^{-1}$ via \eqref{eq:key_identity}). Consequently, a model trained exclusively with \eqref{eq:L_dir} lacks alignment with the NPBS operator structure.

To resolve this misalignment, an intermediate approach views the composition $\mathcal{M}_\theta G_\eta$ as an approximation of $A^{-1}$. Adopting this within the standard residual training paradigm \eqref{eq:L_dir} yields
\begin{equation}
\mathcal{L}_{\mathrm{bs}}^{\ell_2}(\theta)
=
\E_{\bm r\sim\mathcal{D}}
\frac{
\left\lVert A\mathcal{M}_\theta(G_\eta \bm r)-\bm r\right\rVert_{2}
}{
\left\lVert \bm r\right\rVert_{2}
}.
\label{eq:L_bs_l2}
\end{equation}
Although \eqref{eq:L_bs_l2} successfully evaluates the network on the preconditioned inputs $G_\eta \bm r$, it measures the residual in the original Euclidean geometry associated with $A\bm u=\bm f$. Thus, it fails to capture the residual geometry of the left-preconditioned system $G_\eta A \bm u = G_\eta \bm f$. Measuring the residual after left preconditioning, and incorporating the identity \eqref{eq:key_identity}, we propose the metric-matched objective:
\begin{equation}
\mathcal{L}_{\mathrm{bs}}^{R_\eta}(\theta)
=
\E_{\bm r\sim\mathcal{D}}
\frac{
\left\|
(I-G_\eta V_\eta)\,\mathcal{M}_\theta(G_\eta \bm r)-G_\eta \bm r
\right\|_2
}{
\|G_\eta \bm r\|_2
}.
\label{eq:L_bs_Reta_integral}
\end{equation}
In \eqref{eq:L_bs_Reta_integral}, we represent the preconditioned operator as $I-G_\eta V_\eta$ rather than $G_\eta A$. This computationally efficient form requires only a single transform application, mirroring the practical implementation of the NPBS iteration \eqref{eq:npbs_integral}. Unlike \eqref{eq:L_bs_l2}, this proposed objective thoroughly aligns both the input representation and the residual metric with the underlying preconditioned system.




\begin{proposition}[Riesz-map equivalence]
\label{prop:riesz_map}
To make the geometry underlying \eqref{eq:L_bs_Reta_integral} explicit, define the metric induced by the reference Green operator:
\begin{equation}
R_\eta:=G_\eta^\ast G_\eta,\qquad
\|\bm x\|_{R_\eta}:=\|G_\eta \bm x\|_2,\qquad
\langle \bm x,\bm y\rangle_{R_\eta}:=\langle \bm x,R_\eta \bm y\rangle_2.
\end{equation}
Then \eqref{eq:L_bs_Reta_integral} is equivalently
\begin{equation}
\mathcal{L}_{\mathrm{bs}}^{R_\eta}(\theta)
=
\E_{\bm r\sim\mathcal{D}}
\frac{
\left\lVert A\mathcal{M}_\theta(G_\eta \bm r)-\bm r\right\rVert_{R_\eta}
}{
\left\lVert \bm r\right\rVert_{R_\eta}
}.
\label{eq:L_bs_Reta_riesz}
\end{equation}
\end{proposition}

The proof is given in Appendix~\ref{app:proof_riesz_map}.

Proposition~\ref{prop:riesz_map} makes precise that \eqref{eq:L_bs_Reta_integral} and \eqref{eq:L_bs_Reta_riesz} are the same objective written in two equivalent forms. It also clarifies that $\mathcal{L}_{\mathrm{bs}}^{\ell_2}$ and $\mathcal{L}_{\mathrm{bs}}^{R_\eta}$ use the same NPBS-compatible input $G_\eta\bm r$ and the same residual of the original operator, $A\mathcal{M}_\theta(G_\eta\bm r)-\bm r$, but evaluate that residual in different metrics (Euclidean versus $R_\eta$-induced). Thus, relative to \eqref{eq:L_dir}, the progression is: unpreconditioned probes with Euclidean norm ($\mathcal{L}_{\mathrm{dir}}$), then preconditioned probes with Euclidean norm ($\mathcal{L}_{\mathrm{bs}}^{\ell_2}$), and finally preconditioned probes with the $R_\eta$-metric ($\mathcal{L}_{\mathrm{bs}}^{R_\eta}$). A detailed Fourier-domain interpretation of the metric reweighting is given in Appendix~\ref{app:iterp_precond_geometry}.



% \begin{remark}[Training/inference consistency]
% At inference, the NPBS network input is the Born residual
% \begin{equation}
% \bm r_{\mathrm{bs}}^m = G_\eta(\bm f-A\bm u^m).    
% \end{equation}
% The training inputs in \eqref{eq:L_bs_l2} and \eqref{eq:L_bs_Reta_integral} have the same \emph{form} $G_\eta \bm r$. Exact distribution matching holds if $\mathcal{D}$ is chosen as the empirical distribution of iterative residuals (e.g., rollout replay). To mitigate the overhead of dataset preparation, we enforce form-level alignment and treat residual-replay matching as an optional refinement.
% Implementation follows directly from the definitions above:
% (i) fix $\eta$ and the reference operator $L_\eta$ (hence $G_\eta$),
% (ii) train $\mathcal{M}_\theta$ by sampling $\bm r\sim\mathcal{D}$ and minimizing \eqref{eq:L_bs_Reta_integral} (or \eqref{eq:L_bs_Reta_GA}),
% and (iii) at inference iterate \eqref{eq:npbs_integral} until $\|\bm f-A\bm u^m\|_2/\|\bm f\|_2\le \mathrm{rTol}$.
% \end{remark}

The main text presents derivations and operator identities for the Helmholtz setting; CDR and Newton extensions are given in Appendix~\ref{app:cdr_nl_full}, fast-transform verification in Appendix~\ref{app:fft_verify_all}, the pseudospectral recipe in Appendix~\ref{app:fft_ps}, and architecture details in Appendix~\ref{app:architectures}.


%%%%%%%%%%%%%%%%%%%%%%%%%%%%%%%%%%%%%%%%%%%%%



\section{Experiments}
\label{sec:numerical_results}
We organize the evaluation around the heterogeneous Helmholtz equation. 
(i) we isolate the effects of iteration format and training loss across dataset difficulty and wavenumber regimes. 
(ii) we test whether the resulting gains persist across multiple neural-operator backbones. 
(iii) we benchmark the full NPBS solver against classical CBS and shifted-Laplacian preconditioners, both as stationary iterations and inside FGMRES. 
(iv) we evaluate on convection--diffusion--reaction (CDR) systems and Newton linear systems for nonlinear PDEs to demonstrate generality beyond Helmholtz. 
All methods share the unified stopping criterion $\mathrm{rTol}=\norm{\bm f-A\bm u^m}/\norm{\bm f}$.
The primary learned preconditioner uses an FNO backbone (Appendix~\ref{app:architectures}); dataset splits, training hyperparameters, and compute resources are detailed in Appendices~\ref{app:experimental_setup} and~\ref{app:compute_resources}.

\subsection{Heterogeneous Helmholtz Equation}
\label{sec:exp_heter_helm_eq}
For the heterogeneous Helmholtz equation, we select two datasets from OpenFWI~\cite{deng2022openfwi} with progressive difficulty: CurveVel-A (smooth curved velocity profiles) and CurveFault-B (sharp contrasts and fault structures); see Figure~\ref{fig:openfwi_dataset} for representative samples. Both are discretized on a $256\times256$ grid with 32-cell PML boundaries and $12$ points per wavelength ($\texttt{ppw}$). Full dataset and discretization details are given in Appendix~\ref{app:helmholtz_dataset}.

\subsubsection{Iteration format and loss comparison}
We isolate the effect of the network iteration format and the choice of training loss on solver efficiency. The evaluation compares three configurations under a unified stopping rule ($\text{rTol}=10^{-6}$): (i) Direct iteration trained with $\mathcal{L}_{\mathrm{dir}}$, (ii) NPBS iteration trained with $\mathcal{L}_{\mathrm{bs}}^{\ell_2}$, and (iii) NPBS iteration trained with the metric-matched objective $\mathcal{L}_{\mathrm{bs}}^{R_\eta}$. 
\begin{table}[htbp]
\centering
\caption{Iteration format and loss comparison on OpenFWI datasets. Lower $\texttt{ppw}$ corresponds to higher wavenumber. Ratio: iteration count of NPBS+$\mathcal{L}_{\mathrm{bs}}^{R_\eta}$ vs.\ Direct+$\mathcal{L}_{\mathrm{dir}}$.}
\label{tab:loss_formulation}
% \renewcommand{\arraystretch}{1.2}
\resizebox{\linewidth}{!}{%
\begin{tabular}{llccccccc}
\toprule
\multirow{2}{*}{Dataset} & \multirow{2}{*}{$\texttt{ppw}$} & \multicolumn{2}{c}{Direct + $\mathcal{L}_{\mathrm{dir}}$} & \multicolumn{2}{c}{NPBS + $\mathcal{L}_{\mathrm{bs}}^{\ell_2}$} & \multicolumn{2}{c}{NPBS + $\mathcal{L}_{\mathrm{bs}}^{R_\eta}$} & \multirow{2}{*}{Ratio} \\
\cmidrule(lr){3-4}\cmidrule(lr){5-6}\cmidrule(lr){7-8}
& & Avg. Iters & Avg. Time & Avg. Iters & Avg. Time & Avg. Iters & Avg. Time & \\
\midrule
CurveVel-A & 12 & 491 & 0.499\,s & 466 & 0.470\,s & \textbf{300} & \textbf{0.301\,s} & 1.6 \\
\midrule
\multirow{4}{*}{CurveFault-B}
& 24 &  531 & 0.537\,s &  517 & 0.537\,s & \textbf{458} & \textbf{0.403\,s} & 1.2 \\
& 20 &  616 & 0.642\,s &  539 & 0.550\,s & \textbf{475} & \textbf{0.465\,s} & 1.3 \\
& 16 &  895 & 0.914\,s &  842 & 0.848\,s & \textbf{538} & \textbf{0.544\,s} & 1.6 \\
& 12 & 1051 & 1.030\,s &  859 & 0.838\,s & \textbf{556} & \textbf{0.571\,s} & 1.9 \\
\bottomrule
\end{tabular}%
}
\end{table}

Table~\ref{tab:loss_formulation} reports the effect of iteration scheme and training loss across OpenFWI datasets and wavenumber regimes. 
NPBS~+~$\mathcal{L}_{\mathrm{bs}}^{R_\eta}$ consistently requires the fewest iterations and achieves the lowest runtime; its advantage over Direct~+~$\mathcal{L}_{\mathrm{dir}}$ grows both with stronger heterogeneity, from $1.6\times$ on CurveVel-A to $1.9\times$ on CurveFault-B, and with increasing wavenumber, from $1.2\times$ to $1.9\times$ as $\texttt{ppw}$ decreases from 24 to 12. 
These trends indicate that metric matching is most beneficial in harder Helmholtz regimes, while the runtime gains and the small $0.7\%$ training overhead (Appendix~\ref{app:training_convergence}) corroborate the efficient integral-form design in \S\ref{par:efficient_integral}. 



\subsubsection{Cross-backbone robustness of iteration format and loss}
To verify whether the NPBS advantage is architecture-agnostic, we instantiate all three formulations with four backbones: U-Net~\cite{ronneberger2015u}, MgNO~\cite{he2024mgno}, Encoder--solver based on~\cite{lerer2024multigrid}, and MGCFNN~\cite{xie2025mgcfnn}.  
Implementation details and baseline adaptations are described in Appendix~\ref{app:setup_helmholtz}.

Table~\ref{tab:learning_baselines} shows two robust trends. First, for every backbone, NPBS+$\mathcal{L}_{\mathrm{bs}}^{R_\eta}$ yields the lowest iteration count. Second, solve-time ranking differs from iteration ranking: although MGCFNN attains the fewest iterations, MgNO gives the lowest average solve time, reflecting its lightweight design and compile-friendly implementation.

\begin{table}[htbp]
\centering
\caption{Cross-backbone comparison on CurveFault-B with representative learning-based baselines.}
\label{tab:learning_baselines}
% \renewcommand{\arraystretch}{1.2}
\begin{tabular}{lcccccc}
\toprule
\multirow{2}{*}{Backbone} & \multicolumn{2}{c}{Direct + $\mathcal{L}_{\mathrm{dir}}$} & \multicolumn{2}{c}{NPBS + $\mathcal{L}_{\mathrm{bs}}^{\ell_2}$} & \multicolumn{2}{c}{NPBS + $\mathcal{L}_{\mathrm{bs}}^{R_\eta}$} \\
\cmidrule(lr){2-3}\cmidrule(lr){4-5}\cmidrule(lr){6-7}
& Avg. Iters & Avg. Time & Avg. Iters & Avg. Time & Avg. Iters & Avg. Time \\
\midrule
U-Net & 853 & 0.433\,s & 496 & 0.254\,s & 399 & 0.198\,s \\
MgNO & 342 & 0.268\,s & 141 & 0.123\,s & 123 & \textbf{0.107\,s} \\
Encoder--solver & 468 & 0.562\,s & 404 & 0.455\,s & 321 & 0.350\,s \\
MGCFNN & 69 & 0.336\,s & 66 & 0.334\,s & \textbf{48} & 0.235\,s \\
\bottomrule
\end{tabular}
\end{table}

\subsubsection{Comparison with classical iterative methods}
Following the cross-backbone study in Section~\ref{sec:exp_heter_helm_eq}, we select MgNO as the NPBS backbone for its lowest average solve time among the evaluated neural architectures. We compare this instantiation against classical baselines---CBS~\cite{osnabrugge2016convergent} and shifted-Laplacian multigrid~\cite{erlangga2006novel}(SL)---on CurveFault-B, both as stationary iterations and as preconditioners inside FGMRES~\cite{saad1993flexible} with restart length~$50$.

As shown in Table~\ref{tab:helm_cbs_vs_npbs}, NPBS~(MgNO) as a stationary iteration already reduces the average iteration count by $9.1\times$ relative to CBS. Wrapping CBS or SL in FGMRES roughly halves their iteration counts, but the expensive per-step Krylov construction offsets this gain in wall-clock time. In contrast, FGMRES+NPBS~(MgNO) converges in only $37$ iterations with an average time of $0.078$\,s, achieving state-of-the-art performance.
Figure~\ref{fig:openfwi_infer} depicts a representative CurveFault-B sample with its wavefield solution and convergence profiles.


\begin{table}[htbp]
\centering
\caption{Comparison of classical methods and learned NPBS (MgNO backbone) with NPBS + $\mathcal{L}_{\mathrm{bs}}^{R_\eta}$}
\label{tab:helm_cbs_vs_npbs}
% \renewcommand{\arraystretch}{1.2}
% \resizebox{\linewidth}{!}{%
\begin{tabular}{lcc}
\toprule
Method & Avg. Iters & Avg. Time \\
\midrule
CBS & 1124 & 0.086\,s \\
FGMRES+CBS & 547 & 1.820\,s \\
FGMRES+SL & 781 & 4.312\,s \\
NPBS (MgNO) & 123 & 0.107\,s \\
FGMRES+NPBS (MgNO) & \textbf{37} & \textbf{0.078}\,s \\
\bottomrule
\end{tabular}%
% }
\end{table}


\begin{figure}[htbp]
    \centering
    \includegraphics[width=\linewidth]{openfwi_infer_comp.png}
    \caption{Iteration comparison of different strategies on a representative instance from the OpenFWI CurveFault-B dataset.}
    \label{fig:openfwi_infer}
\end{figure}


\subsection{Algorithmic Generalization Beyond the Helmholtz Equation}
\label{sec:generalization}
To verify that the metric-matched formulation generalizes beyond the Helmholtz equation, we evaluate the same iteration format and loss design on two additional PDE settings: a linear convection--diffusion--reaction (CDR) system and Newton linear systems for a nonlinear PDE. 
Dataset and setup details are in Appendix~\ref{app:setup_generalization} and \ref{app:dataset_figures}.

\paragraph{Linear Convection--Diffusion--Reaction (CDR) Systems}
We consider the linear CDR equation
\begin{equation}
  - \nabla \cdot (\kappa \nabla u) + \mathbf{v}\cdot \nabla u + \sigma u = f
  \quad \text{in } \Omega \subset \mathbb{R}^d,
\end{equation}
posed with periodic boundary conditions. Following the Helmholtz study, we compare three formulations and hold the architecture fixed so that only the iteration format and training objective vary.
We construct a synthetic dataset on a $128\times 128$ periodic grid with spatially heterogeneous coefficients (dataset generation details and representative samples are provided in Appendix~\ref{app:cdr_dataset}).

As shown in Table~\ref{tab:other_pde_solver_cmp} and Figure~\ref{fig:cdr_iter_his}, the advantage of NPBS with $\mathcal{L}_{\mathrm{bs}}^{R_\eta}$ observed on the Helmholtz equation carries over to the CDR setting: the metric-matched formulation consistently achieves the fastest convergence among the three objectives.

\begin{figure}[ht]
    \centering
    \includegraphics[width=0.85\linewidth]{cdr_iter_his.png}
    \caption{Representative CDR test case. Left: solution field $u$. Right: residual histories.}
    \label{fig:cdr_iter_his}
\end{figure}

\paragraph{Newton Linear Systems for Nonlinear PDEs}
We consider the non-convex problem with multiple solutions studied in~\cite{hao2024newton}:
\begin{equation}
\begin{cases}
  -\Delta u(x,y) - u(x,y)^2 = -s \sin(\pi x)\sin(\pi y), & (x,y) \in \Omega,\\[2mm]
  u(x,y) = 0, & (x,y) \in \partial\Omega,
\end{cases}
\end{equation}
with $f(x,y):=-s\sin(\pi x)\sin(\pi y)$ and define nonlinear residual $\mathcal{F}_{\mathrm{nl}}(u):=-\Delta u-u^2-f$. 
At Newton iterate $u^{(m)}$, the linearized correction $\delta u^{(m)}$ solves
\begin{equation}\label{eq:newton}
\begin{cases}
  -\Delta (\delta u^{(m)}) - 2u^{(m)}\,\delta u^{(m)} = -\mathcal{F}_{\mathrm{nl}}(u^{(m)}), & (x,y) \in \Omega,\\[1mm]
  \delta u^{(m)} = 0, & (x,y) \in \partial\Omega,
\end{cases}
\end{equation}
followed by the update $u^{(m+1)}=u^{(m)}+\delta u^{(m)}$. 
In every Newton step, the learned preconditioner is applied to the inner linear system~\eqref{eq:newton}.
Following~\cite{hao2024newton}, discretization and dataset details are in Appendix~\ref{app:newton_dataset}. For this Newton benchmark, the inner linear solves use $\mathrm{rTol}=10^{-4}$.

Table~\ref{tab:other_pde_solver_cmp} shows that the ranking observed in the linear benchmarks carries over to the Newton inner solves: the metric-matched NPBS requires the fewest iterations.
Figure~\ref{fig:newton_iter} further shows that the learned preconditioner integrates stably within the outer Newton loop.
\begin{figure}[htbp]
    \centering
    \includegraphics[width=0.95\linewidth]{newton_iter.png}
    \caption{FNO preconditioner with NPBS+$\mathcal{L}_{\mathrm{bs}}^{R_\eta}$ embedded in Newton iterations.}
    \label{fig:newton_iter}
\end{figure}

\begin{table}[htbp]
\centering
\renewcommand{\arraystretch}{1.2}
\caption{Iteration format and loss comparison across other PDEs.}
\label{tab:other_pde_solver_cmp}
\begin{tabular}{l cc cc}
\toprule
\multirow{2}{*}{\textbf{Method}} & \multicolumn{2}{c}{\textbf{CDR Benchmark}} & \multicolumn{2}{c}{\textbf{Newton Inner-Solve}} \\
\cmidrule(lr){2-3} \cmidrule(lr){4-5}
& Avg.\ Iters & Avg.\ Rel Res & Avg.\ Iters & Avg.\ Rel Res \\
\midrule
Direct + $\mathcal{L}_{\mathrm{dir}}$          & 391          & $9.84\times10^{-7}$ & 25          & $8.98\times10^{-5}$ \\
NPBS + $\mathcal{L}_{\mathrm{bs}}^{\ell_2}$    & 167          & $9.57\times10^{-7}$ & 11          & $7.20\times10^{-5}$ \\
NPBS + $\mathcal{L}_{\mathrm{bs}}^{R_\eta}$    & \textbf{96}  & $8.68\times10^{-7}$ & \textbf{9}  & $6.75\times10^{-5}$ \\
\bottomrule
\end{tabular}
\end{table}


\section{Conclusion}
\label{sec:conclusion}
We introduced Neural Preconditioned Born Series (NPBS), a learned iterative preconditioning framework that operates in Born-preconditioned residual coordinates. The identity $I-G_\eta V_\eta = G_\eta A$ shows that these coordinates are exactly those of shifted-Laplacian left preconditioning, and the induced metric $R_\eta=G_\eta^\ast G_\eta$ leads to a metric-matched objective $\mathcal{L}_{\mathrm{bs}}^{R_\eta}$ that aligns training with inference.
Experiments on heterogeneous Helmholtz problems show that metric-matched NPBS consistently reduces iteration counts, with larger gains in more difficult regimes and robustness across neural backbones. When embedded in FGMRES, NPBS achieves the lowest average solve time among the evaluated methods. Results on CDR systems and Newton linear systems for nonlinear PDEs further indicate that residual-metric matching is useful beyond Helmholtz. Future work includes lighter neural preconditioners, sharper rules for choosing the reference operator and damping parameter, and extensions to larger-scale three-dimensional and time-dependent problems.

%%%%%%%%%%%%%%%%%%%%%%%%%%%%%%%%%%%%%%%%%%%%%%%%%%%%%%%%%%%%
\bibliographystyle{plainnat}
\bibliography{references}


\appendix

\section{Proofs and Supporting Analysis}
\label{app:proofs}

\subsection{Proof of Proposition~\ref{prop:equivalence_spectral}}
\label{app:proof_equivalence_spectral}

\begin{proof}

\emph{Identity.}\; From $G_\eta L_\eta=I$ and $A=L_\eta-V_\eta$:
\[
I-G_\eta V_\eta
=
G_\eta L_\eta - G_\eta V_\eta
=
G_\eta(L_\eta - V_\eta)
=
G_\eta A.
\]
\emph{Spectral bounds.}\; Since $\rho(G_\eta V_\eta)\le\|G_\eta V_\eta\|_2\le q$, every eigenvalue $\mu$ of $G_\eta V_\eta$ satisfies $|\mu|\le q$. Eigenvalues of $I-G_\eta V_\eta$ are $1-\mu$, giving \eqref{eq:disk_bound}. The condition $\|G_\eta V_\eta\|_2<1$ ensures the Neumann series converges:
\begin{equation}
(I-G_\eta V_\eta)^{-1}=\sum_{j=0}^{\infty}(G_\eta V_\eta)^j,
\label{eq:neumann_bound}
\end{equation}
so that
\[
\|I-G_\eta V_\eta\|_2\le 1+q,\qquad
\|(I-G_\eta V_\eta)^{-1}\|_2\le \sum_{j\ge 0}q^j=\frac{1}{1-q},
\]
which yields
\begin{equation}
\kappa_2(I-G_\eta V_\eta)\le \frac{1+q}{1-q}.
\end{equation}
\end{proof}
Part of the conclusion was also noted in~\cite{huang2021finite}.
\subsection{Proof of Proposition~\ref{prop:riesz_map}}
\label{app:proof_riesz_map}

\begin{proof}
Using \eqref{eq:key_identity}, \eqref{eq:L_bs_Reta_integral} can be rewritten as
\begin{equation}
\mathcal{L}_{\mathrm{bs}}^{R_\eta}(\theta)
=
\E_{\bm r\sim\mathcal{D}}
\frac{
\left\|
G_\eta A\,\mathcal{M}_\theta(G_\eta \bm r)-G_\eta \bm r
\right\|_2
}{
\|G_\eta \bm r\|_2
}.
\label{eq:L_bs_Reta_GA}
\end{equation}
By definition, $\|\bm x\|_{R_\eta}=\|G_\eta \bm x\|_2$. Substituting
$\bm x = A\mathcal{M}_\theta(G_\eta \bm r)-\bm r$ gives
$\|A\mathcal{M}_\theta(G_\eta \bm r)-\bm r\|_{R_\eta}
=
\|G_\eta A\mathcal{M}_\theta(G_\eta \bm r)-G_\eta \bm r\|_2$,
and similarly $\|\bm r\|_{R_\eta}=\|G_\eta \bm r\|_2$, yielding \eqref{eq:L_bs_Reta_riesz}.
\end{proof}

\subsection{Interpretation in preconditioned geometry.}
\label{app:iterp_precond_geometry}
Relative to \eqref{eq:L_dir}, \eqref{eq:L_bs_l2} and \eqref{eq:L_bs_Reta_integral} replace the direct regression target $A^{-1}$ with the preconditioned target $(G_\eta A)^{-1}$ by feeding $G_\eta \bm r$ into $\mathcal{M}_\theta$. The key difference lies in the residual metric: $\mathcal{L}_{\mathrm{bs}}^{\ell_2}$ measures the physical residual $e(\bm r;\theta):=A\,\mathcal{M}_\theta(G_\eta \bm r)-\bm r$ in the Euclidean norm $\|e(\bm r;\theta)\|_2$, whereas $\mathcal{L}_{\mathrm{bs}}^{R_\eta}$ measures it in the preconditioned norm $\|e(\bm r;\theta)\|_{R_\eta}=\|G_\eta e(\bm r;\theta)\|_2$. By Proposition~\ref{prop:riesz_map}, the latter is exactly the residual ratio evaluated in the same $R_\eta$-geometry induced by the preconditioned coordinates used by NPBS inference. For the constant-coefficient Helmholtz reference, the Fourier symbol
\begin{equation}
\widehat{G_\eta}(\bm \xi)=\frac{1}{|\bm \xi|^2-k_0^2-i\eta}
\end{equation}
shows explicitly how the metric reweights residual modes in an $\eta$-controlled way. In this sense, optimization is formulated in the left-preconditioned geometry associated with $G_\eta A=I-G_\eta V_\eta$. When $\|G_\eta V_\eta\|<1$ (Proposition~\ref{prop:equivalence_spectral}), the spectrum of $G_\eta A$ is concentrated near $1$, so different modal components are brought to a more comparable scale. This provides an interpretable explanation for why the metric-matched objective may exhibit more favorable optimization behavior in practice than training directly against $A$.

\subsection{Role of \texorpdfstring{$\eta$}{eta} in Metric and Iteration Coordinates}
\label{app:role_eta}

The shift $\eta$ defines both the left preconditioner $G_\eta:=L_\eta^{-1}$ and the induced residual metric
$R_\eta:=G_\eta^\ast G_\eta$.
Consequently, the $R_\eta$-weighted residual norm satisfies
\begin{equation}
\|\bm f-A\bm u\|_{R_\eta}^2 := (\bm f-A\bm u)^\ast R_\eta (\bm f-A\bm u)=\|G_\eta(\bm f-A\bm u)\|_2^2,
\end{equation}
so $\eta$ controls the geometry in which residuals are measured and corrected, not only numerical damping.
In practice, we tune $\eta$ via a small validation sweep and then keep it fixed between training and inference
for each dataset to preserve objective--inference consistency.
Heuristically, increasing $\eta$ can move the reference inverse farther from resonance and may improve the contraction behavior of the fixed-point map, whereas excessively large damping can also weaken the usefulness of the resulting correction directions.

\section{Full Derivations for CDR and Nonlinear Extensions}
\label{app:cdr_nl_full}
This appendix gives the full technical derivations deferred from the main text.

\subsection{Convection-Diffusion-Reaction (Linear)}
Consider
\begin{equation}
\mathcal{L}_{\mathrm{cdr}}u
=
-\nabla\cdot(\kappa(x)\nabla u)
\;+\;
\mathbf{v}(x)\cdot\nabla u
\;+\;
\sigma(x)u
=
f.
\label{eq:cdr_pde}
\end{equation}
Choose background constants $\kappa_0,\mathbf{v}_0,\sigma_0$ and define
\begin{equation}
\kappa=\kappa_0+\delta\kappa,\qquad
\mathbf{v}=\mathbf{v}_0+\delta\mathbf{v},\qquad
\sigma=\sigma_0+\delta\sigma.
\end{equation}
Set
\begin{equation}
\mathcal{L}_{0,\mathrm{cdr}}u
:=
-\kappa_0\Delta u
\;+\;
\mathbf{v}_0\cdot\nabla u
\;+\;
\sigma_0 u,
\label{eq:cdr_L0}
\end{equation}
and define $\mathcal{V}_{\mathrm{cdr}}$ by $\mathcal{L}_{\mathrm{cdr}}=\mathcal{L}_{0,\mathrm{cdr}}-\mathcal{V}_{\mathrm{cdr}}$, i.e.,
\begin{equation}
\mathcal{V}_{\mathrm{cdr}}u
=
\nabla\cdot(\delta\kappa\nabla u)
\;-\;
\delta\mathbf{v}\cdot\nabla u
\;-\;
\delta\sigma\,u.
\label{eq:cdr_V}
\end{equation}
Then
\begin{equation}
\mathcal{L}_{0,\mathrm{cdr}}u
=
\mathcal{V}_{\mathrm{cdr}}u+f.
\end{equation}
With $G_{\mathrm{cdr}}:=\mathcal{L}_{0,\mathrm{cdr}}^{-1}$:
\begin{equation}
(I-G_{\mathrm{cdr}}\mathcal{V}_{\mathrm{cdr}})u
=
G_{\mathrm{cdr}}f
=
G_{\mathrm{cdr}}\mathcal{L}_{\mathrm{cdr}}u.
\label{eq:cdr_integral}
\end{equation}
In discrete form ($A_{\mathrm{cdr}},V_{\mathrm{cdr}},G_{\mathrm{cdr}}$):
\begin{equation}
I-G_{\mathrm{cdr}}V_{\mathrm{cdr}}
=
G_{\mathrm{cdr}}A_{\mathrm{cdr}}.
\label{eq:cdr_identity}
\end{equation}
Therefore, the CDR-specific instance of the metric-matched Born Series inspired objective is
\begin{equation}
\mathcal{L}_{\mathrm{bs}}^{R_{\mathrm{cdr}}}(\theta)
=
\E_{r\sim\mathcal{D}}
\frac{
\norm{
(I-G_{\mathrm{cdr}}V_{\mathrm{cdr}})
\mathcal{M}_\theta(G_{\mathrm{cdr}}r)
-G_{\mathrm{cdr}}r
}
}{
\norm{G_{\mathrm{cdr}}r}
}
\label{eq:cdr_bs_loss}
\end{equation}
or equivalently
\begin{equation}
\mathcal{L}_{\mathrm{bs}}^{R_{\mathrm{cdr}}}(\theta)
=
\E_{r\sim\mathcal{D}}
\frac{
\norm{
G_{\mathrm{cdr}}A_{\mathrm{cdr}}
\mathcal{M}_\theta(G_{\mathrm{cdr}}r)
-G_{\mathrm{cdr}}r
}
}{
\norm{G_{\mathrm{cdr}}r}
}.
\label{eq:cdr_bs_loss_eq}
\end{equation}
with $R_{\mathrm{cdr}}:=G_{\mathrm{cdr}}^\ast G_{\mathrm{cdr}}$. This is exactly the same template as $\mathcal{L}_{\mathrm{bs}}^{R_\eta}$ in Section~\ref{sec:bs_loss}, under the substitution $(G_\eta,A,V_\eta)\mapsto(G_{\mathrm{cdr}},A_{\mathrm{cdr}},V_{\mathrm{cdr}})$.

For periodic domains, $G_{\mathrm{cdr}}$ is a convolution operator with Fourier symbol
\begin{equation}
\widehat{G}_{\mathrm{cdr}}(\xi)
=
\frac{1}{\kappa_0|\xi|^2+i\mathbf{v}_0\cdot\xi+\sigma_0},
\label{eq:cdr_G_symbol}
\end{equation}
so high-frequency diffusion-dominated modes are attenuated in training. As in Helmholtz, optimizing against $G_{\mathrm{cdr}}A_{\mathrm{cdr}}$ instead of raw $A_{\mathrm{cdr}}$ improves conditioning and better matches the residual distribution in preconditioned iterations.

\subsection{Nonlinear PDE via Newton-Born Linearization}
For the nonlinear example
\begin{equation}
\mathcal{F}_{\mathrm{nl}}(u):=-\Delta u-u^2-f=0,
\label{eq:nonlinear_residual}
\end{equation}
Newton at iterate $u^{(m)}$ solves
\begin{equation}
J_m\,\delta u^{(m)} = r^{(m)},
\qquad
J_m:=-\Delta-2u^{(m)},
\qquad
r^{(m)}:=-\mathcal{F}_{\mathrm{nl}}(u^{(m)}),
\label{eq:newton_linear_system}
\end{equation}
followed by $u^{(m+1)}=u^{(m)}+\delta u^{(m)}$.

\paragraph{Discrete realization used in this work (Dirichlet + FDM + DST).}
For the Newton experiments in Section~\ref{sec:numerical_results}, we use homogeneous Dirichlet boundary conditions on $\Omega=[0,1]^2$ and a five-point finite-difference discretization on interior unknowns. With $h=1/(N+1)$, define the discrete negative Laplacian
\begin{equation}
L_{\mathrm{D}}u
=
\frac{1}{h^2}
\begin{bmatrix}
0 & -1 & 0\\
-1 & 4 & -1\\
0 & -1 & 0
\end{bmatrix}*u,
\end{equation}
so the discrete Jacobian is
\begin{equation}
J_m^{\mathrm{disc}} = L_{\mathrm{D}} - 2\,\mathrm{diag}(u^{(m)}).
\end{equation}
In our implementation, the $-2u^{(m)}$ term is replaced by its spatial average $-2\bar{u}^{(m)}$ in the preconditioning pipeline, where $\bar{u}^{(m)}$ denotes the spatial mean of $u^{(m)}$.

To build a Born-style preconditioned form for each Newton step, pick a reference Jacobian
\begin{equation}
J_{0,m}:=-\Delta-2\bar{u}^{(m)}+\alpha_m \I,
\qquad
\bar{u}^{(m)}:=\text{spatial average of }u^{(m)},
\qquad
\alpha_m\ge 0,
\label{eq:newton_ref}
\end{equation}
whose discrete counterpart is $J_{0,m}^{\mathrm{disc}}=L_{\mathrm{D}}-2\bar u^{(m)}I+\alpha_m I$ (or $J_{0,m}^{\mathrm{disc}}\approx L_{\mathrm{D}}+\alpha_m I$ in the Laplacian-dominant implementation). Equivalently, if one denotes the code-level Laplace operator by $\Delta_h$, then $L_{\mathrm{D}}=-\Delta_h$ under the sign convention used here. Define
\begin{equation}
V_m:=J_{0,m}-J_m.
\label{eq:newton_vm}
\end{equation}
Then $J_m=J_{0,m}-V_m$, $G_m:=J_{0,m}^{-1}$, and
\begin{equation}
(I-G_mV_m)\,\delta u^{(m)}=G_m r^{(m)}
\quad\Leftrightarrow\quad
I-G_mV_m=G_mJ_m.
\label{eq:newton_born_identity}
\end{equation}
Hence the Newton-state instance of the metric-matched Born Series objective is
\begin{equation}
\mathcal{L}_{\mathrm{bs}}^{R_{\mathrm{N}}}(\theta)
=
\E_{(u^{(m)},r)\sim\mathcal{D}_{\mathrm{N}}}
\frac{
\norm{
(I-G_mV_m)\mathcal{M}_\theta(G_m r)-G_m r
}
}{
\norm{G_m r}
},
\label{eq:nl_bs_loss}
\end{equation}
equivalently
\begin{equation}
\mathcal{L}_{\mathrm{bs}}^{R_{\mathrm{N}}}(\theta)
=
\E_{(u^{(m)},r)\sim\mathcal{D}_{\mathrm{N}}}
\frac{
\norm{
G_mJ_m\mathcal{M}_\theta(G_m r)-G_m r
}
}{
\norm{G_m r}
}.
\label{eq:nl_bs_loss_eq}
\end{equation}
where the per-state metric is $R_m:=G_m^\ast G_m$ (so $R_{\mathrm{N}}$ denotes the Newton-trajectory family $\{R_m\}$). This is the same metric-matched objective family as $\mathcal{L}_{\mathrm{bs}}^{R_\eta}$, with state-dependent references.

\paragraph{Discussion (nonlinear).}
The loss is computed on \emph{linearized} systems along Newton trajectories, not directly on the nonlinear map. This has two effects:
\begin{enumerate}
\item It preserves the same preconditioned-coordinate advantage as the linear case (better-conditioned operator equation).
\item It aligns training with actual inference, where the solver repeatedly tackles Jacobian systems with changing $u^{(m)}$.
\end{enumerate}
In practice, one may use an inner iteration
\begin{equation}
\delta u_{k+1}^{(m)}
=
\delta u_k^{(m)}
\;+\;
\mathcal{M}_\theta\!\left(
G_m\left(r^{(m)}-J_m\delta u_k^{(m)}\right)
\right),
\label{eq:newton_inner_iter}
\end{equation}
then set $\delta u^{(m)}\approx \delta u_{K}^{(m)}$ and update $u^{(m+1)}=u^{(m)}+\delta u^{(m)}$.

\section{Verification: Fast Transform Implementation of Green Operators}
\label{app:fft_verify_all}
We verify, case by case, when Green-operator application is fast. A concrete Fourier pseudospectral discretization recipe for $L_\eta$ and $G_\eta$ is provided in Appendix~\ref{app:fft_ps}.

\begin{proposition}[Transform-diagonalizable Green operators]
\label{prop:fft_all_cases}
On a uniform grid, if the chosen reference operator has constant coefficients and compatible boundary conditions (periodic, or transform-compatible Dirichlet/Neumann), then its discrete matrix is diagonalizable by a fast trigonometric transform. Consequently, applying the corresponding discrete Green operator costs $\mathcal{O}(N\log N)$.
\end{proposition}

\begin{proof}[Constructive verification]
Let $L_0$ be any constant-coefficient reference operator and $G_0=L_0^{-1}$. Under periodic boundary conditions, the discrete $L_0$ is block-circulant with circulant blocks, hence
\[
L_0 = F^{-1}\Lambda F,\qquad
G_0 = F^{-1}\Lambda^{-1}F,
\]
where $F$ is the multidimensional DFT matrix and $\Lambda$ contains the symbol values. Therefore
\[
G_0 q = \mathcal{F}^{-1}\!\left(\frac{\widehat{q}}{\lambda(\xi)}\right),
\]
implemented by FFT/IFFT in $\mathcal{O}(N\log N)$.
For Dirichlet or Neumann boundaries, the same statement holds with sine/cosine transforms (DST/DCT), which are FFT-based and have the same complexity order.
\end{proof}

\paragraph{Helmholtz case.}
With $L_\eta=-\Delta-(k_0^2+i\eta)$, the symbol is
\[
\lambda_{\mathrm{H}}(\xi)=|\xi|^2-(k_0^2+i\eta),
\]
so
\[
G_\eta q = \mathcal{F}^{-1}\!\left(\frac{\widehat{q}}{\lambda_{\mathrm{H}}(\xi)}\right).
\]
Hence the Born residual $G_\eta r$ is FFT-fast under the stated boundary setting.

\paragraph{CDR case.}
With $\mathcal{L}_{0,\mathrm{cdr}}=-\kappa_0\Delta+\mathbf{v}_0\cdot\nabla+\sigma_0$, the symbol is exactly \eqref{eq:cdr_G_symbol} denominator:
\[
\lambda_{\mathrm{cdr}}(\xi)=\kappa_0|\xi|^2+i\mathbf{v}_0\cdot\xi+\sigma_0.
\]
Thus
\[
G_{\mathrm{cdr}} q = \mathcal{F}^{-1}\!\left(\frac{\widehat{q}}{\lambda_{\mathrm{cdr}}(\xi)}\right),
\]
again $\mathcal{O}(N\log N)$.

\paragraph{Newton linear systems for nonlinear PDEs.}
For the Newton experiments, we use homogeneous Dirichlet boundaries and five-point FDM on interior nodes. The reference Jacobian is diagonalized by the 2D DST-I (not periodic FFT). Writing
\[
J_{0,m}^{\mathrm{disc}} = S^{-1}\Lambda_{\mathrm{N},m}S,
\]
with $S$ the 2D DST-I transform, the preconditioned application is
\[
G_m q = S^{-1}\!\left(\frac{Sq}{\lambda_{p,q}^{(m)}}\right),
\]
where for the Laplacian-dominant implementation
\[
\lambda_{p,q}^{(m)}
\approx
\frac{4}{h^2}\left[\sin^2\!\left(\frac{p\pi h}{2}\right)+\sin^2\!\left(\frac{q\pi h}{2}\right)\right]
\;(+\,\alpha_m\ \text{if used}).
\]
If the mean-reaction shift is retained, the full eigenvalues of $J_{0,m}^{\mathrm{disc}}$ are the Laplacian eigenvalues above shifted by $\alpha_m-2\bar u^{(m)}$. Therefore the Newton reference inverse is applied by DST/IDST in $\mathcal{O}(N\log N)$.

\paragraph{Boundary-condition note for all cases.}
In scattering settings, we use an absorbing layer (e.g., sponge/PML-style damping) on a finite box and choose a \emph{constant-coefficient} reference operator on that computational domain; this is the operator inverted by FFT in the Born-inspired loss/iteration. If strict Dirichlet/Neumann boundaries are required, use DST/DCT (or FFT-based embedding), preserving near-identical complexity.
In particular, the experiments on Newton linear systems for nonlinear PDEs use homogeneous Dirichlet boundaries and DST-based inversion of the FDM Laplacian reference.

\section{Fourier Pseudospectral Discretization of \texorpdfstring{$L_\eta$ and $G_\eta$}{L eta and G eta}}
\label{app:fft_ps}
This appendix gives the concrete discretization used when $L_\eta$ and $G_\eta=L_\eta^{-1}$ are implemented by FFT under periodic boundary conditions. In numerical-analysis terminology, this is a \emph{Fourier pseudospectral} (or Fourier collocation) discretization: derivatives are represented by exact Fourier multipliers, while coefficient multiplications are applied pointwise in physical space~\cite{boyd2001chebyshev,trefethen2000spectral,shen2011spectral}.

\paragraph{Periodic grid and Fourier modes.}
Let $\Omega=\prod_{\alpha=1}^d [0,L_\alpha)$ and use a uniform grid with $N_\alpha$ points in direction $\alpha$:
\[
x_{j_\alpha}=j_\alpha h_\alpha,\qquad
h_\alpha=L_\alpha/N_\alpha,\qquad
j_\alpha=0,\dots,N_\alpha-1.
\]
For grid field $u_j=u(x_j)$, define the DFT coefficients $\widehat u_n$ with multi-index
\[
n=(n_1,\dots,n_d),\qquad
n_\alpha\in\Big\{-\frac{N_\alpha}{2},\dots,\frac{N_\alpha}{2}-1\Big\},
\]
and physical wavenumbers
\[
\xi_{\alpha,n}=\frac{2\pi n_\alpha}{L_\alpha},\qquad
|\xi_n|^2=\sum_{\alpha=1}^d \xi_{\alpha,n}^2.
\]

\paragraph{Discrete symbol of $L_\eta$.}
For
\[
L_\eta=-\Delta-(k_0^2+i\eta),\qquad \eta>0,
\]
the Fourier symbol is
\[
\lambda_n = |\xi_n|^2-(k_0^2+i\eta),
\]
so
\[
\widehat{L_\eta u}_n=\lambda_n\,\widehat u_n.
\]
Hence $L_\eta$ is diagonal in Fourier space, and its inverse acts modewise:
\[
\widehat{(G_\eta q)}_n = \frac{\widehat q_n}{\lambda_n}.
\]
Because $\eta>0$, $\lambda_n$ has nonzero imaginary part and the modal inversion is well-defined.

\paragraph{FFT implementation of $G_\eta q$.}
Given $q$, compute:
\begin{enumerate}
\item $\widehat q=\mathcal{F}(q)$ (FFT),
\item $\widehat z_n=\widehat q_n/\lambda_n$ (pointwise complex division),
\item $z=\mathcal{F}^{-1}(\widehat z)$ (IFFT).
\end{enumerate}
Then $z=G_\eta q$. The complexity is $\mathcal{O}(N\log N)$ with $N=\prod_\alpha N_\alpha$.

\paragraph{Relation to $A$ and $V_\eta$.}
For heterogeneous Helmholtz,
\[
A u = -\Delta u - k(x)^2u,\qquad
V_\eta u=(k(x)^2-k_0^2-i\eta)u.
\]
In Fourier pseudospectral form:
\[
-\Delta u = \mathcal{F}^{-1}\!\big(|\xi|^2\widehat u\big),
\]
while $k(x)^2u$ and $V_\eta u$ are computed pointwise in physical space. This derivative-in-Fourier / product-in-physical split is exactly the pseudospectral pattern~\cite{boyd2001chebyshev,trefethen2000spectral}. For strongly non-smooth coefficients, standard de-aliasing (e.g., $2/3$ truncation) can be used~\cite{shen2011spectral}.

\paragraph{Boundary-condition scope.}
The FFT diagonalization above is exact for periodic boundaries. For homogeneous Dirichlet/Neumann boundaries, the same constant-coefficient reference operators are diagonalized by sine/cosine transforms (DST/DCT), with the same $\mathcal{O}(N\log N)$ complexity class. In this manuscript, $G_\eta$ in CBS/MgNO-BS is implemented with the periodic/transform-compatible constant-coefficient reference operator to preserve fast application.



\section{Architectural Instantiations}
\label{app:architectures}
The main text is not architecture-centered; NPBS is compatible with different neural operators. For reproducibility, we briefly record the primary FNO instantiation used in most experiments.

\paragraph{FNO instantiation and design principle.}
Across Helmholtz, CDR, and Newton linear systems for nonlinear PDEs, we use an FNO-based preconditioner $\mathcal{M}_\theta$. Let $a$ denote PDE coefficient fields (e.g., $v$ in Helmholtz, or $a=(\kappa,\mathbf{v},\sigma)$ in CDR), and let $r$ denote the iterative residual. The design explicitly applies a \emph{nonlinear} transform to coefficient features and a \emph{linear} transform to residual features, then fuses both branches:
\begin{equation}
\mathcal{M}_\theta(r;a)
=
\mathcal{P} \Big(
\mathrm{LFL} \Big[
\mathrm{NLFL}(\Phi_\theta(a))
\odot
\mathrm{LFL}(\Psi_\theta(r))
\Big]
\Big).
\label{eq:fno_impl}
\end{equation}
The lifting/project mappings are
\begin{equation}
\Phi_\theta(a) = W_a * \mathrm{LayerNorm}(\mathrm{pad}(a)),
\qquad
\Psi_\theta(r) = W_r * \mathrm{pad}(r),
\label{eq:fno_lift}
\end{equation}
and
\begin{equation}
\mathcal{P}(z) = W_p * \mathrm{unpad}(z).
\label{eq:fno_proj}
\end{equation}
Here $\mathrm{NLFL}(\Phi_\theta(a))$ is used for nonlinear coefficient encoding, while the residual branch keeps linear processing through $\mathrm{LFL}(\Psi_\theta(r))$, ensuring linearity with respect to $r$.

The Fourier blocks are
\begin{equation}
\mathrm{LFL}(x)
=
\mathcal{F}^{-1}\!\big(\mathcal{W}^{\mathrm{F}}_{\theta}\odot \mathcal{F}(x)\big)
\;+
\mathcal{W}^{\mathrm{loc}}_{\theta} * x,
\label{eq:lfl}
\end{equation}
\begin{equation}
\mathrm{NLFL}(x)
=
\varphi\!\big(\mathrm{LFL}(x)\big),
\label{eq:nlfl}
\end{equation}
where $\mathcal{F}$ and $\mathcal{F}^{-1}$ are Fourier transforms, $\mathcal{W}^{\mathrm{F}}_{\theta}$ are truncated-mode spectral weights, $\mathcal{W}^{\mathrm{loc}}_{\theta}$ is a local $1\times1$ convolution, and $\varphi$ is GELU. In implementation, coefficient features pass through stacked NLFL blocks, while the residual branch uses a single LFL block.

\section{Dataset Illustration}
\label{app:dataset_figures}

\subsection{Heterogeneous Helmholtz Equation}
\label{app:helmholtz_dataset}

We use velocity models from the OpenFWI benchmark~\cite{deng2022openfwi}. For both CurveVel-A and CurveFault-B, the training set contains $3{,}000$ samples and the test set contains $50$ samples. The computational domain is discretized on a $256\times256$ grid with a 32-cell perfectly matched layer (PML) boundary, yielding a padded grid of size $288\times288$. Spatial velocity fields range from approximately $1500$ to $4500$\,m/s, and the grid resolution ensures at least $12$ points per wavelength. Representative velocity models are shown in Figure~\ref{fig:openfwi_dataset}.

\begin{figure}[ht]
    \centering
    \includegraphics[width=0.95\linewidth]{openfwi_curvevel_A_and_curvefault_B.png}
    \caption{Representative heterogeneous velocity models from the OpenFWI datasets. Top row: CurveVel-A, featuring relatively smooth, curved velocity transitions. Bottom row: CurveFault-B, characterized by sharp discontinuities, strong contrasts, and complex fault structures.}
    \label{fig:openfwi_dataset}
\end{figure}

\subsection{CDR Dataset Generation}
\label{app:cdr_dataset}

The CDR benchmark uses a $128\times128$ periodic grid. The training set contains $4000$ samples and the test set contains $100$ samples. The diffusion field $\kappa$ is drawn from a Gaussian random field (GRF); the velocity $\mathbf{v}$ is obtained from a GRF stream function $\psi$ via
\begin{equation}
v_x=\frac{\partial \psi}{\partial y},\qquad v_y=-\frac{\partial \psi}{\partial x},
\end{equation}
so that incompressibility $\nabla \cdot \mathbf{v}=0$ is enforced by construction; the source term $f$ is a Gaussian bump with randomly sampled center and width. Representative samples are shown in Figure~\ref{fig:cdr_data_samples}.

\begin{figure}[ht]
    \centering
    \includegraphics[width=0.95\linewidth]{cdr_data_samples.png}
    \caption{Representative samples from the CDR benchmark. Each row shows the diffusion field $\kappa$, the incompressible velocity field $\mathbf{v}$ visualized by streamlines, and the source term $f$.}
    \label{fig:cdr_data_samples}
\end{figure}

\subsection{Newton Benchmark Data}
\label{app:newton_dataset}

Following~\cite{hao2024newton}, the domain $\Omega=[0,1]^2$ is discretized with $N=64$ interior grid points (grid spacing $h=1/65$) and source amplitude $s=1600$. Training data consists of $5000$ Newton trajectory pairs, and the test set contains $100$ samples. Supervision is obtained by generating Newton trajectories with sparse direct solves. To ensure that the trajectories probe distinct solution branches, initial guesses are constructed by perturbing a reference solution with low-frequency sine-series noise,
\begin{equation}
u^0=u^\star+\alpha\sum_{i,j=1}^{4}c_{ij}\sin(i\pi x)\sin(j\pi y),\quad
c_{ij}\sim\mathcal{N}(0,0.1^2),\ \alpha=4.0.
\end{equation}
Representative branches obtained in this way are shown in Figure~\ref{fig:newton_diff_sols}.

\begin{figure}[ht]
    \centering
    \includegraphics[width=0.65\linewidth]{newton_diff_sols.png}
    \caption{Multiple solutions of the nonlinear PDE under different initial guesses.}
    \label{fig:newton_diff_sols}
\end{figure}

%%%%%%%%%%%%%%%%%%%%%%%%%%%%%%%%%%%%%%%%%%%%%%%%%%%%%%%%%%%%


\section{Experimental Setup}
\label{app:experimental_setup}

This appendix details the dataset splits, backbone architectures, and training hyperparameters for all experiments reported in Section~\ref{sec:numerical_results}.

\subsection{Heterogeneous Helmholtz Equation}
\label{app:setup_helmholtz}

\paragraph{Dataset.}
We use velocity models from the OpenFWI benchmark~\cite{deng2022openfwi}. For both CurveVel-A and CurveFault-B, the training set contains $3{,}000$ samples and the test set contains $50$ samples. The computational domain is discretized on a $256\times256$ grid with a 32-cell perfectly matched layer (PML) boundary, yielding a padded grid of size $288\times288$ (after accounting for the absorbing boundary on each side).

\paragraph{FNO backbone (primary).}
The primary preconditioner follows the FNO instantiation of Eq.~\eqref{eq:fno_impl} with $3$ NLFL layers, $25$ hidden channels, and $45$ retained Fourier modes per layer. The residual branch uses a single LFL with full-bandwidth modes ($N_{\mathrm{pad}}/2+1$). Training uses the Adam optimizer for $400$ epochs with batch size $30$, initial learning rate $10^{-3}$, and a step-decay schedule (factor $0.8$ every $20$ epochs). Inference is performed in float64 with a relative tolerance of $10^{-6}$ and a maximum of $2{,}500$ iterations.

\paragraph{Comparison backbones.}
To verify that the NPBS gains are architecture-agnostic, we evaluate additional backbones under the same NPBS+$\mathcal{L}_{\mathrm{bs}}^{R_\eta}$ framework. All comparison models are trained for $400$ epochs with batch size $30$, learning rate $10^{-3}$, and step-decay (factor $0.8$ every $20$ epochs). Their architecture-specific parameters are:
\begin{itemize}
    \item \textbf{MgNO}~\cite{he2024mgno}: A multigrid neural operator with $32$ channels, $4$-level V-cycle structure with $[2,2]$ pre-/post-smoothing iterations per level, and GELU activation.
    \item \textbf{MGCFNN}~\cite{xie2025mgcfnn}: A multigrid convolutional Fourier neural network with $32$ channels, $3$-level hierarchy. Level~0 uses spatial $5\times5$ convolutions (1 sweep); levels~1 and~2 use complex Fourier layers with wavenumber cutoffs $15$ and $10$ respectively (2 and 4 sweeps).
    \item \textbf{Encoder-Solver}~\cite{lerer2024multigrid}: A U-Net-style encoder-decoder with $32$ channels, $4$ levels, and implicit weight-sharing across decoder iterations.
    \item \textbf{U-Net}~\cite{ronneberger2015u}: A standard $4$-level U-Net with $32$ channels and skip connections; the rhs-path convolutions are bias-free to preserve linearity in the residual.
    \item \textbf{DeepONet}~\cite{lu2021learning}: A DeepONet variant with $32$ basis functions, $64$ trunk hidden channels, and $4$ blocks.
\end{itemize}
For runtime, we use \texttt{torch.compile} when stable. MGCFNN is evaluated without compilation because \texttt{torch.compile} led to divergence during inference. The original Encoder--solver is supervised and does not preserve the linear structure required in our iterative setting, which leads to divergence; we therefore adopt a linearized variant that retains its core multigrid-inspired components.

\subsection{Algorithmic Generalization Beyond the Helmholtz Equation}
\label{app:setup_generalization}

\subsubsection{Convection--Diffusion--Reaction (CDR)}
\label{app:setup_cdr}

\paragraph{Dataset.}
The CDR benchmark uses a $128\times128$ periodic grid. The training set contains $4{,}000$ samples and the test set contains $400$ samples. Coefficients are generated as described in Appendix~\ref{app:cdr_dataset}.

\paragraph{FNO backbone.}
The preconditioner uses the same FNO design principle as in the Helmholtz setting (Eq.~\eqref{eq:fno_impl}), with $25$ hidden channels, $3$ NLFL layers with $32$ retained Fourier modes per layer, and zero-padding of $8$ cells. Training uses Adam for $60$ epochs with batch size $32$, initial learning rate $10^{-3}$, and step-decay (factor $0.8$ every $3$ epochs). Inference uses a relative tolerance of $10^{-6}$.

\subsubsection{Newton Linear Systems for Nonlinear PDEs}
\label{app:setup_newton}

\paragraph{Dataset.}
Following~\cite{hao2024newton}, the domain $\Omega=[0,1]^2$ is discretized with $N=64$ interior grid points (grid spacing $h=1/65$) and the source amplitude is $s=1600$. Training data consists of $10{,}000$ Newton trajectory pairs (combining two trajectory sets of $5{,}000$ each from different solution branches), and the test set contains $1{,}000$ samples.

\paragraph{FNO backbone.}
The preconditioner uses $25$ hidden channels, $2$ NLFL layers with $20$ retained Fourier modes, and zero-padding of $8$ cells. Training uses Adam for $200$ epochs with batch size $50$, initial learning rate $10^{-4}$, and weight decay $10^{-6}$. The inner linear solver uses a relative tolerance of $10^{-4}$.

%%%%%%%%%%%%%%%%%%%%%%%%%%%%%%%%%%%%%%%%%%%%%%%%%%%%%%%%%%%%

\section{Compute Resources}
\label{app:compute_resources}

All experiments were conducted using PyTorch 2.6.0 with CUDA 12.4 on a single workstation equipped with an Intel Xeon Platinum 8468 CPU and one NVIDIA H100 GPU.
All solvers---including classical baselines such as CBS and the shifted-Laplacian preconditioning---are implemented entirely in PyTorch and accelerated via CUDA, ensuring that wall-clock comparisons reflect algorithmic differences rather than implementation-level disparities.
No distributed training, large-scale pretraining, or fine-tuning of pretrained foundation models was involved; the main computational cost comes from training lightweight neural preconditioners and iterative solver inference.

%%%%%%%%%%%%%%%%%%%%%%%%%%%%%%%%%%%%%%%%%%%%%%%%%%%%%%%%%%%%

\section{Training Convergence}
\label{app:training_convergence}

Figure~\ref{fig:train_val_loss_curvefault_b} shows the training and validation rTol histories for the three loss configurations on CurveFault-B. 
%Notably, the ordering of monitored rTol during training is opposite to the inference iteration counts reported in Table~\ref{tab:loss_formulation}, confirming that lower probe-level error does not directly translate to faster iterative convergence.

\begin{figure}[ht]
    \centering
    \includegraphics[width=0.75\linewidth]{train_val_loss.png}
    \caption{Training and validation rTol histories on CurveFault-B.}
    \label{fig:train_val_loss_curvefault_b}
\end{figure}

\begin{table}[htbp]
\centering
\caption{Training statistics on CurveFault-B under the same FNO setup.}
\label{tab:curvefault_b_train_stats}
\renewcommand{\arraystretch}{1.15}
\begin{tabular}{lccc}
\toprule
Method & Train rTol & Val rTol & Train time \\
\midrule
Direct + $\mathcal{L}_{\mathrm{dir}}$ & $8.91\times 10^{-2}$ & $8.77\times 10^{-2}$ & 52~min~18~s \\
NPBS + $\mathcal{L}_{\mathrm{bs}}^{\ell_2}$ & $8.99\times 10^{-2}$ & $8.90\times 10^{-2}$ & 52~min~35~s \\
NPBS + $\mathcal{L}_{\mathrm{bs}}^{R_\eta}$ & $9.20\times 10^{-2}$ & $9.04\times 10^{-2}$ & 52~min~39~s \\
\bottomrule
\end{tabular}
\end{table}


%%%%%%%%%%%%%%%%%%%%%%%%%%%%%%%%%%%%%%%%%%%%%%%%%%%%%%%%%%%%
\section{Limitations}
\label{app:limitations}

\paragraph{Dimensionality and grid structure.}
Our numerical evaluation focuses on two-dimensional problems on uniform structured grids. We do not consider three-dimensional settings, where data generation, memory use, and neural-preconditioner application may be substantially more expensive. 
We also do not evaluate unstructured meshes or complex geometries.
\paragraph{Theoretical guarantees.}
The spectral clustering and condition-number bounds in Proposition~\ref{prop:equivalence_spectral} rely on $\|G_\eta V_\eta\|_2\le q<1$. For regimes that violate this condition, we do not provide a theoretical convergence or conditioning guarantee.

\paragraph{Applicability to PDE classes.}
NPBS is not suitable for arbitrary PDEs. It relies on a reference--perturbation split $A=L_0-V$ with an invertible reference operator $L_0$, so that the Born-preconditioned form $I-G_0V$ with $G_0=L_0^{-1}$ is well defined. This structure is natural for the Helmholtz, CDR, and Newton linear systems studied here, where a reference operator captures the main part of the PDE and the remaining coefficient or state dependence is treated as a perturbation. PDEs for which no meaningful or stable reference--perturbation split is available fall outside the present framework.

\newpage
% \input{checklist.tex}


\end{document}
